\phantomsection\addcontentsline{toc}{chapter}{Abstract} % Don't remove me!
\chapter*{\centering{\Large Abstract}}

Gaze dynamics during visual search are rich both in space and time, and their study has taught us much about computation across the brain. The process of active vision during search is a feedback loop of gathering new information based on what has been learnt, involving large areas of the brain at once. Gaze trajectories are of unique interest in being a direct and easily acquired measurement of neural attention. Current work on visual search is divided in focus and object of study, with tension between biological and predictive accuracy, as well as a gap between studying the large spatiotemporal scales of gaze through jumps in focus called saccades, and the small scales between these saccades called fixations. In this thesis, we unify these interests and scales under a single model which treats gaze as a complex dynamical system, as well as a sophisticated probabilistic computational system. Using the framework of fractional neural sampling (FNS) we focus this model through the key dynamical signatures found in visual gaze dynamics, global spatial superdiffusion and the nested oscillations of regular saccades which modulate fine spatial sampling within fixations. This model can thus be used to understand the range of brain function across subjects, including as a non-invasive biomarker to understand disorders of attention and perception (ADHD, schizophrenia, autism). As a sampling model, this same work can inform generalised probabilistic computation strategies. This forms the foundation for this model as providing a direct measurement of neural attention to understand the computational function of the brain across all levels.